\section{Dataset}
\label{sec:dataset}

The benchmark combines three corpora. Each one was built by its own single-axis benchmark
using the same recipe. We wrote families of code generators, sampled tasks from each
family, and measured several runs of each task at different input sizes in Kubernetes
pods. \Cref{tab:corpora} summarizes the three corpora, and \cref{app:families} lists every
family. All of these measurements describe the workload itself, independent of how its
pod was submitted. \Cref{sec:twins} then works out what a Job pod and a Flyte pod each
need, starting from the same measurement.

\begin{table}[tbp]
\centering\small
\caption{The three corpora. RightSizeGPU's 2,052 unusable runs: 1,569 censored (out of
memory on the largest shape they reached), 51 failed, and 432 the combined bench drops
because their peak filled the device, which makes the measured peak only a lower bound.
RightSizeMem's 166 unusable runs exceeded its 20\,GiB guard or 900\,s timeout.}
\label{tab:corpora}
\resizebox{\linewidth}{!}{%
\begin{tabular}{@{}llrrrrrrl@{}}
\toprule
corpus & release & families & held out & tasks/fam. & runs/task & measured & usable & labels \\
\midrule
RightSizeMem & v2    & 12 & 2 & 400 & 3 & 14,400 & 14,234 & peak memory; mean cores \\
RightSizeCPU & big-1 & 11 & 2 & 450 & 2 &  9,900 &  9,900 & peak memory; busiest 15\,s \\
RightSizeGPU & big-2 & 10 & 3 & 500 & 2 & 10,000 &  7,948 & peak GPU memory, sharding \\
\midrule
combined     & k8s-2 & 33 & 7 &     &   & 34,300 & 32,082 & \\
\bottomrule
\end{tabular}}
\end{table}

\subsection{Families, tasks and runs}
A \emph{family} is a Python generator. Given a seed, it writes a task function and the
specification of its inputs, and it draws a few hidden settings, which we call knobs,
that decide how input bytes turn into resource use. A pandas family, for example,
decides whether to project columns, which dtype backend to use and whether to downcast
floats. A training family picks width, depth, batch size, mixed precision and
checkpointing, and a CPU family picks a thread-pool size or leaves it unset. The knobs
show up only in the code an approach reads. They are never given to an approach directly,
and we keep them only to slice the results.

Each run scales the leading dimension of the task's inputs by a factor drawn
log-uniformly from $[0.25, 4]$ for RightSizeMem and from $[0.5, 2]$ for the CPU and GPU
corpora. Runs are shuffled, so a task's history can be larger or smaller than the current
run. The inputs are synthetic: normal floats, uniform integers, dictionary-encoded
categories, random strings, Zipf-distributed word ids for text, and random bytes for
images, written as Parquet, CSV, JSONL or \code{.npy}. CPU and GPU workloads loop until a
per-task time budget runs out (30--60\,s for CPU and 45--90\,s for GPU), while RightSizeMem
workloads run to completion.

\subsection{RightSizeMem v2: peak memory}
RightSizeMem has twelve families, chosen because they turn input bytes into memory in
different ways. They include pandas group-bys and joins, polars pipelines in eager, lazy
and streaming modes, DuckDB queries with and without a memory limit, chunked CSV and
JSONL readers, NumPy linear algebra with optional memory mapping, scikit-learn
estimators, TF-IDF, a small PyTorch model on CPU, image batches and pure-Python grouping.

Each run uses three processes. A child process writes the inputs to a temporary
directory, a fresh interpreter runs the task body on those files, and the parent reads
the child's \code{ru\_maxrss} and CPU time from \code{wait4}. BLAS and OpenMP are limited
to one thread, and workloads set their own threads in code. The workload's peak memory
$w$ is its \code{ru\_maxrss}. As we show in \cref{sec:validation}, on Linux this value
includes the measuring parent's own peak at the time of \code{exec}, about 158\,MiB here,
which makes about a third of the smaller workloads' labels one cell higher than needed.

\subsection{RightSizeCPU big-1: the busiest 15 seconds}
RightSizeCPU has eleven families, chosen for how they use cores, since CPU need depends on
parallelism much more than on input size. They cover NumPy BLAS with and without a thread
cap, polars and DuckDB with explicit thread settings, scikit-learn ensembles, serial
pandas, Python thread pools that hold or release the GIL, process pools, a CPU PyTorch
model, pyarrow CSV parsing, an I/O-bound pool and a task with two phases.

CPU time for the whole process tree is sampled every second. Each pod runs two
measurements side by side on disjoint sets of seven cores, and a check of 44 runs
against single-run pods put all 44 within 25\% of each other. The CPU label is the
busiest 15-second average, because that is the resolution at which cluster monitoring
(Prometheus scraping cAdvisor every 15\,s) sees a container. Peak memory is measured as in
RightSizeMem.

\subsection{RightSizeGPU big-2: peak GPU memory across a machine ladder}
RightSizeGPU has ten PyTorch families: a tabular MLP, a CNN, a CNN limited by its data
loader, a GPT language model, a transformer encoder, a recommender with sharded
embeddings, blocked matrix multiplication, nearest-neighbour search, fine-tuning of
decoders from 0.6B to 7B parameters split across GPUs, and generation with a KV cache.
NVML reports memory and utilization on every visible GPU every half second. Each run
starts on a T4 and moves up the machine ladder (\cref{tab:ladder}) whenever it runs out
of device memory, and a run that runs out of memory on the largest machine it reached is
censored. The oracle machine for a run is the cheapest one it fits on. For machines a run
was not measured on, we model the fit from its per-GPU peak and, for sharded code, its
imbalance across GPUs.

\subsection{Splits and history}
Runs are split by task, never by run. A hash of the task id sends 70\% of each family's
tasks to training, 10\% to validation and 20\% to test. The seven held-out families go
entirely to a separate \emph{unseen} split. The test split therefore measures how well an
approach handles new tasks from familiar families, and the unseen split measures how well
it handles new kinds of code. Every setting that an approach tunes is chosen on
validation, which contains familiar families only.

The $k$-th run of a task carries the task's earlier runs as its history. A task's first
run has no history at all, and these cold starts make up 33--50\% of each split.

\begin{table}[tbp]
\centering\small
\caption{Usable runs per split. Each run appears twice in the bench, once as a Job and
once as a Flyte task, in the same split.}
\label{tab:splits}
\begin{tabular}{@{}lrrrrl@{}}
\toprule
source & train & val & test & unseen & unseen families \\
\midrule
mem & 8,457 & 1,138 & 2,316 & 2,323 & duckdb\_sql, image\_batch \\
cpu & 5,628 &   838 & 1,634 & 1,800 & duckdb\_sql, pyarrow\_csv \\
gpu & 3,968 &   542 & 1,145 & 2,293 & recsys\_train, knn\_search, llm\_generate \\
\midrule
runs        & 18,053 & 2,518 &  5,095 &  6,416 & \\
pods (rows) & 36,106 & 5,036 & 10,190 & 12,832 & \\
\bottomrule
\end{tabular}
\end{table}
